\section{Data}\label{sec:data}

The analysis uses the CDFI Fund's NMTC Public Data Release covering
FY2003--FY2022 award years (released June 2024), which reports 19{,}907
QLICI transactions across 8{,}024 projects with origination years
2001--2022. Transactions are aggregated to the project level; the project
is the unit at which total cost---and therefore leverage---is defined.
Cleaning steps, field definitions, and file hashes are documented in the
public repository's data dictionary and provenance files, and the full
pipeline from the raw workbook to every table in this paper runs from a
single script sequence.

The outcome is project-level leverage,
\[
\mathrm{Leverage}_i \;=\; \frac{\mathrm{ProjectCost}_i}{\mathrm{QLICI}_i},
\]
which I refer to as leverage throughout. Leverage is bounded below at 1 by
construction, and the mass at 1 collects projects whose recorded cost equals
the subsidized investment alone.

A word is owed on what this ratio measures, since the blended-finance
framing invites a reading the data do not support. The QLICI is the
investment the intermediary places, and it is not itself a federal outlay.
The government's cost is the 39\% credit claimed against the qualified
equity investment that funds the QLICI, spread across seven years, so a
dollar of QLICI corresponds to roughly thirty-nine cents of tax expenditure
under the substantially-all requirement that ties the two magnitudes
together. The aggregates make the distinction concrete: \$120.9 billion of
project cost against \$66.6 billion of QLICI leaves \$54.3 billion of
other capital, which is \$0.82 for every dollar of subsidized investment
and about \$2.09 for every dollar of implied tax expenditure. I therefore
read leverage as a measure of financing structure. None of the paper's
identification turns on the choice of denominator, since a constant
rescaling leaves the rural coefficient's sign, its relative magnitude, and
every test reported below unchanged. The baseline outcome is
winsorized to $[1, 20]$ to discipline a long right tail of large
capital-stack deals; Section~\ref{sec:robust} shows the headline result is
insensitive to the cap, including its removal.

Table~\ref{tab:desc} reports descriptives. Non-metro projects are 19.5\% of
the sample---again, the statutory line---and lag metro projects in both
mean (1.73 versus 1.99) and median (1.07 versus 1.19) leverage.
Figure~\ref{fig:dist} shows the shape of the gap: non-metro mass piles up
near the 1.0 floor while the metro distribution carries a fatter right
shoulder.

\begin{table}[t]
\centering
\caption{NMTC deployment by metro status, FY2001--FY2022.}
\label{tab:desc}
\begin{tabular}{lrr}
\toprule
 & Metro & Non-metro \\
\midrule
Projects & 6,463 & 1,561 \\
QLICI deployed (\$M) & 53,566 & 13,045 \\
Total project cost (\$M) & 97,151 & 23,751 \\
Mean leverage & 1.99 & 1.73 \\
Median leverage & 1.19 & 1.07 \\
\bottomrule
\end{tabular}

\end{table}

\begin{figure}[t]
\centering
\includegraphics[width=0.48\textwidth]{1_allocation_timeseries.png}\hfill
\includegraphics[width=0.48\textwidth]{2_non_metro_share_timeseries.png}
\caption{Deployment over time. Left: annual QLICI dollars, metro and
non-metro. Right: the non-metro dollar share against the 20\% statutory
line.}
\label{fig:deploy}
\end{figure}

\begin{figure}[t]
\centering
\includegraphics[width=0.85\textwidth]{3_leverage_distribution.png}
\caption{Project-level leverage distributions by metro status. Non-metro
mass concentrates at the $1.0\times$ floor (zero private mobilization);
the metro distribution has the fatter right shoulder.}
\label{fig:dist}
\end{figure}
